\textbf{Маргинальные распределения}

\begin{definition}
Пусть $P$ -- распределение на $(\mathbb{R}^n, B(\mathbb{R}^n)), i \in \{1, \ldots, n\}, B \in B(\mathbb{R})$
\end{definition}

Тогда $P_i(B) = P(\R^{i-1} \times B \times \R^{n-i})$ -- распределение на $(\R, B(\R))$. Оно называется \textit{маргинальным}. (Идейно -- что-то вроде проекции на оси)

Довольно просто проверить корректность определения (что это действительно распределение).

По маргинальным распределениям можно построить многомерное (не единственным образом, однако, каким-то можно).

Наиболее естественный такой:
 
Пусть $P_1, \ldots, P_n$ -- распределения на $(\R, B(\R))$. Положим $P(B_1 \times \ldots \times B_n) = P_1(B_1) \cdot \ldots \cdot P_n(B_n)$ -- \textit{декартово произведение распределений} $P_1, \ldots, P_n$. У него $F(x_1, \ldots, x_n) = F_1(x_1) \cdots \ldots \cdot F_n(x_n)$.
 
\begin{lemmanote} 
Это распределение достаточно задать на декартовых произведениях борелевских множеств, дальше оно однозначно продолжится на всю сигма-алгебру. Более того, на самом деле, нам даже достаточно задать на декартовых произведениях лучей (потому что, как мы знаем, лучами можно задать всю борелевскую сигма-алгебру на $\R$)\\
\end{lemmanote}


\textbf{Дискретные распределения}

Тут все аналогично одномерному случаю (примеры можно получить из одномерного с помощью декартова произведения).

$X \subset \R^n$ -- не более, чем счетное, $P(X) = 1 \Rightarrow P$ -- \textit{дискретное}

Теперь введем \textit{дискретную плотность} $p: X \to [0, 1], p(x) = P(\{x\})$

$\forall B \in B(\R^n) : P(B) = \sum\limits_{x \in X \cap B}p(x)$

\begin{lemmanote}
как и в одномерном случае, дискретное распределение можно задать, задав его значение в каждой точке. Оно будет единственным образом восстанавливаться из такого описания.\\
\end{lemmanote}

\textbf{Абсолютно непрерывное распределение}

\begin{definition}
Если $p: \R^n \to \R_+$ такова, что $\forall x \in \R^n$
\end{definition}

$F(x) = \int\limits_{-\infty}^{x_1}\ldots\int\limits_{-\infty}^{x_n} p(t_1, \ldots, t_n)dt_1\ldots dt_n$, то $P$ -- абсолютно непрерывное распределение с плотностью $p$.

\textbf{Свойства}

\begin{enumerate}
    \item $\forall B \in B(\R^n) : P(B) = \int\limits_B p dt_1\ldots dt_n$
    \item Если $p : \R^n \to \R_+, \int\limits_{\R^n}p dt_1 \ldots dt_n = 1$, то $p$ -- плотность некоторого распределения вероятностей
    \item Если $F$ дифференцируема по каждому аргументу, то $P$ -- абсолютно непрерывное распределение с плотностью $p(t_1, \ldots, t_n) = \frac{\partial^n F}{\partial t_1 \ldots \partial t_n}$
\end{enumerate}

Свойства 1) и 2) доказываются аналогично одномерной ситуации. Свойство 3) доказывается проверкой того, что интеграл плотности равен $F$, что, очевидно, верно.

\begin{statement}
Если $P_1, \ldots, P_n$ -- абсолютно непрерывные распределения на $(\R, B(\R))$ с плотностями $p_1, \ldots, p_n$, то $P_1 \times \ldots \times P_n$ -- абсолютно непрерывное распределение на $(\R^n, B(\R^n))$ с плотностью $p_1(t_1) \cdot \ldots \cdot p_n(t_n)$
\end{statement}

\begin{proof}
Положим $p(t_1, \ldots, t_n) = p_1(t_1) \cdot \ldots \cdot p_n(t_n)$. Проверим, что ее интеграл равен $F$:

$\int\limits_{-\infty}^{x_1}\ldots\int\limits_{-\infty}^{x_n} p(t_1, \ldots, t_n)dt_1\ldots dt_n = \int\limits_{-\infty}^{x_1}\ldots\int\limits_{-\infty}^{x_n}  p_1(t_1) \cdot \ldots \cdot p_n(t_n)dt_1\ldots dt_n = \int\limits_{-\infty}^{x_1}p_1(t_1)dt_1\ldots\int\limits_{-\infty}^{x_n}p_n(t_n)dt_n = F_1(x_1)\cdot \ldots \cdot F_n(x_n)$
\end{proof}

\begin{lemmanote}
дискретное распределение абсолютно непрерывно относительно считающей меры (в точках множества $M$ мера равна $1$, в остальных -- $0$)
\end{lemmanote}